\section{犯罪形态 - 犯罪中止}
\begin{enumerate}
  \item 成立要件。
        \begin{enumerate}
          \item 前提条件。\bs{能不能继续犯罪}。
                \begin{enumerate}
                  \item 纯粹主观说。完全看罪犯主观。\al{强奸太丑放弃。主观认为不能继续犯罪 \yjt 被迫放弃 \yjt 未遂}。
                  \item 纯粹客观说。从纯粹的物理、自然角度去判断。\al{强奸对象是亲妹妹 \yjt 能 \yjt 主动放弃 \yjt 中止}。
                  \item \zd{客观说}。社会一般人角度。\al{强奸对象是亲妹妹 \yjt 不能 \yjt 被迫放弃 \yjt 未遂}。
                \end{enumerate}
          \item 主观条件。想不想。主动放弃。
          \item 结论。\zd{能继续而放弃，中止}。常考。
                \begin{enumerate}
                  \item 害怕处罚。\al{盗窃害怕当场/日后被抓 \yjt 未遂/中止}。
                  \item 遇到熟人。\al{拦路抢劫碰到普通/亲人 \yjt 中止/未遂}。
                  \item 客观认识错误。
                        \begin{enumerate}
                          \item \al{盗窃以为是主人回来了，其实是邻居脚步声 \yjt 主观怎么想就怎么定}。
                          \item 通过欺骗让犯罪人产生认识错误。\al{强奸骗有性病，甲被吓跑了 \yjt 主观认为被迫放弃 \yjt 未遂}。
                        \end{enumerate}
                  \item 特定对象障碍。\zd{特定对象不存在，一律未遂}。
                \end{enumerate}
          \item 有效性。
                \begin{enumerate}
                  \item 防止措施。可能、实际的有效性 \yjt 没有 \yjt（特殊）未遂。没考过
                  \item 特殊的中止。因果关系，介入因素。
                \end{enumerate}
        \end{enumerate}
  \item 分类。
        \begin{enumerate}
          \item 未实行终了的中止。自动放弃即可。\al{拿剑要捅却没捅}。
          \item 实行终了的中止。
                \begin{enumerate}
                  \item 自动放弃。
                  \item 实际有效防止的措施，否则拟定未遂。例外。\zd{是否存在介入因素可能阻断了这种实际有效性}。
                  \item 考试。犯罪行为 \yjt 中止 \yjt 抢救措施 \yjt 介入因素 \yjt 实害结果。
                \end{enumerate}
        \end{enumerate}
  \item 处罚。\zd{有/没有造成损害结果 \yjt 应当减轻/免除}。\bs{注意判断这个结果是不是因造成的}。
\end{enumerate}
